\documentclass[12pt]{article}
\usepackage[T1]{fontenc}
\usepackage[utf8]{inputenc}
\usepackage{lmodern}
\usepackage{textcomp}
\usepackage{enumitem}
\usepackage{geometry}
\usepackage{graphicx}
\usepackage{amsmath, amssymb}
\geometry{margin=1in}
\begin{document}
\begin{center}
Physics - Graduate Level Assessment 3 \
Total Marks: 40 \
Answer all questions.
\end{center}

\section*{Multiple Choice (10 marks)}
\begin{enumerate}[label=\arabic*.]
\item When an increase in speed doubles the momentum of a moving body its kinetic energy
\begin{enumerate}[label=(\alph*)]
    \item decreases slightly
    \item increases but does not reach double
    \item remains the same
    \item halves
    \item quadruples
\end{enumerate}
\item The flight of a blimp best illustrates
\begin{enumerate}[label=(\alph*)]
    \item Archimedes' principle
    \item Pascal's principle
    \item Newton's first law
    \item The law of conservation of energy
    \item Hooke's law
\end{enumerate}
\item How do astronomers think Jupiter generates its internal heat?
\begin{enumerate}[label=(\alph*)]
    \item the gravitational pull from its moons
    \item heat produced by its magnetic field
    \item by contracting changing gravitational potential energy into thermal energy
    \item heat absorption from the Sun
    \item nuclear fusion in the core
\end{enumerate}
\item An ultrasonic transducer used for medical diagnosis oscillates at 6.7 Mhz.How long does
each oscillation take, and what is the angular frequency? (Unit: 10\^{}7 rad/s)
\begin{enumerate}[label=(\alph*)]
    \item 5.3
    \item 1.5
    \item 4.2
    \item 8.4
    \item 3.8
\end{enumerate}
\item Why are there no impact craters on the surface of Io?
\begin{enumerate}[label=(\alph*)]
    \item Any craters that existed have been eroded through the strong winds on Io's surface.
    \item Io is protected by a layer of gas that prevents any objects from reaching its surface.
    \item The gravitational pull of nearby moons prevent objects from hitting Io.
    \item Io's magnetic field deflects incoming asteroids, preventing impact craters.
    \item Io did have impact craters but they have all been buried in lava flows.
\end{enumerate}
\end{enumerate}

\section*{True / False (10 marks)}
\textit{Answer True or False}
\begin{enumerate}[label=\arabic*.]
\setcounter{enumi}{5}
\item True or False: An aircraft can produce a sonic boom while flying at a constant speed below the speed of sound.
\item True or False: The angular acceleration of the turntable is 0.58 sec\^{}-2, indicating a slower increase to 33.3 rpm in the same time interval.
\item True or False: The binding energy of the 4He nucleus is approximately 27.5 MeV, suggesting a weaker nuclear interaction than observed.
\item True or False: A particle moving at v=0.60c will decay after traveling only 100 m in the lab frame due to its increased speed.
\item True or False: Mars' axial tilt is less pronounced than Earth's, contributing to its more extreme seasonal variations in the southern hemisphere.
\end{enumerate}

\section*{Long Form Response (20 marks)}
\textit{Answer each question with a clear and complete explanation.}
\begin{enumerate}[label=\arabic*.]
\setcounter{enumi}{10}
\item Discuss the principles of thin-film interference and calculate the required thickness of
an MgF₂ coating on a glass surface to minimize reflection at a wavelength of 550 nm.
\item Discuss the implications of relativistic length contraction as two spaceships approach
Earth at equal speeds from opposite directions, given that a meterstick on one is measured
to be 60 cm by an occupant of the other. What is the speed of each spaceship relative to
the observer on Earth, and how does this scenario illustrate the principles of special
relativity?
\end{enumerate}

\vfill
\noindent\textit{End of Paper}
\end{document}